\documentclass[11pt,a4paper]{ctexart}
\usepackage[a4paper,margin=1in]{geometry}
\usepackage{booktabs}
\usepackage{longtable}
\usepackage{array}
\usepackage{hyperref}
\usepackage{enumitem}
\usepackage{xcolor}

\hypersetup{
  colorlinks=true,
  linkcolor=blue,
  urlcolor=blue,
  pdftitle={GameAgent Voyager 运行报告},
  pdfauthor={OpenCode}
}

\setlist[itemize]{leftmargin=1.8em}
\setlist[enumerate]{leftmargin=2.2em}

\title{GameAgent: Voyager Minecraft 运行与代码修复报告}
\author{基于 MineDojo/Voyager 的工程化分支}
\date{2026-04-29}

\begin{document}

\maketitle

\begin{abstract}
本报告总结了当前 GameAgent 仓库中 Voyager 分支的真实运行结果、关键代码改动、稳定性修复，以及本次实验结束时能够从日志与 checkpoint 中确认的终止原因。该分支的主要目标是在真实 Minecraft 1.19 LAN 世界中持续运行 Voyager，并通过 OpenAI-compatible 接口接入 \texttt{opencode.ai/zen/go/v1}。截至当前快照，系统已经完成 27 个任务，学得 28 个技能，进度推进到铁镐、铁靴与基础洞穴资源阶段。
\end{abstract}

\section{项目背景与目标}

本仓库基于上游 \href{https://github.com/MineDojo/Voyager}{MineDojo/Voyager}，但目标与原始演示仓库不同，重点在于：

\begin{itemize}
  \item 让 Voyager 在真实 Minecraft LAN 世界中连续运行，而不是仅作为论文代码示例。
  \item 使用 OpenAI-compatible API，而非固定依赖 OpenAI 官方端点。
  \item 保留断点续跑、长时间日志监控、以及运行稳定性修复。
  \item 输出可交付的工程文档与运行结果报告。
\end{itemize}

本次使用的核心入口脚本为 \texttt{run\_voyager.py}，通过 \texttt{.env.local} 读取环境变量，目标模型为 \texttt{kimi-k2.6}，API Base 为 \texttt{https://opencode.ai/zen/go/v1}。

\section{关键代码修改}

表 \ref{tab:files} 总结了本分支中最重要的工程性改动。

\begin{longtable}{p{0.28\textwidth}p{0.64\textwidth}}
\caption{关键修改文件与作用}\label{tab:files}\\
\toprule
文件 & 作用 \\
\midrule
\endfirsthead
\toprule
文件 & 作用 \\
\midrule
\endhead
\texttt{run\_voyager.py} & 统一加载 \texttt{.env.local}，清理代理环境变量，配置 OpenAI-compatible 端点，并将环境步进超时提高到 300 秒。 \\
\texttt{run.py} & 去除硬编码密钥，改为从环境变量读取配置，作为短轮次 smoke test 入口。 \\
\texttt{voyager/utils/llm\_utils.py} & 为 LLM 调用增加应用层重试和外层硬超时，避免模型调用无异常但长期卡死。 \\
\texttt{voyager/utils/fake\_embeddings.py} & 当使用自定义 OpenAI-compatible 接口时启用 \texttt{FakeEmbeddings}，避免 embedding 兼容性问题。 \\
\texttt{voyager/agents/action.py} & Action agent 改用统一的重试/超时逻辑。 \\
\texttt{voyager/agents/curriculum.py} & Curriculum、QA、decomposition 全部改用统一 LLM 调用包装。 \\
\texttt{voyager/agents/critic.py} & Critic 调用统一到应用层重试逻辑。 \\
\texttt{voyager/agents/skill.py} & Skill 描述生成统一到应用层重试逻辑。 \\
\texttt{voyager/env/bridge.py} & 增强 \texttt{/step} 与 \texttt{/start} 的恢复能力，mineflayer 重启后会进行重试，不再因为短暂竞态直接终止整个学习流程。 \\
\texttt{voyager/env/mineflayer/index.js} & 修复 bot 生命周期空指针问题，修复 stuck recovery 中对空 block 列表错误取值，增强 \texttt{/step} 阶段的局部 bot 捕获。 \\
\texttt{voyager/control\_primitives/mineBlock.js} & 增加掉落等待和二次收集逻辑，减轻“方块已挖但掉落尚未入包”的情况。 \\
\bottomrule
\end{longtable}

\section{运行结果}

截至当前 checkpoint，系统已经完成下列量化结果：

\begin{itemize}
  \item 已完成任务：27 个
  \item 已失败任务：3 个
  \item 已学得技能：28 个
\end{itemize}

\subsection{已完成任务}

\begin{enumerate}
  \item Mine 1 wood log
  \item Mine 3 spruce\_log
  \item Craft 4 spruce\_planks
  \item Craft 1 crafting\_table
  \item Craft 8 spruce\_planks
  \item Craft 1 wooden\_pickaxe
  \item Mine 3 cobblestone
  \item Mine 1 coal\_ore
  \item Craft 4 sticks
  \item Craft 1 stone\_pickaxe
  \item Mine 1 copper\_ore
  \item Mine 8 cobblestone
  \item Craft 1 furnace
  \item Smelt 4 raw\_copper
  \item Mine 1 lapis\_ore
  \item Mine 4 coal\_ore
  \item Mine 1 iron\_ore
  \item Craft 1 stone\_sword
  \item Equip stone\_sword
  \item Mine 2 iron\_ore
  \item Smelt 3 raw\_iron
  \item Craft 1 iron pickaxe
  \item Mine 3 iron\_ore
  \item Mine 4 iron\_ore
  \item Smelt 4 raw\_iron
  \item Craft 1 iron\_boots
  \item Equip iron\_boots
\end{enumerate}

\subsection{失败任务}

\begin{itemize}
  \item Craft 12 spruce\_planks
  \item Mine 1 spruce\_leaves
  \item Kill 1 skeleton
\end{itemize}

\subsection{技能库结果}

当前已学得技能数为 28，覆盖了木制工具、石制工具、熔炉、煤/铜/铁/青金石挖掘、以及铁靴装备等流程。代表性技能包括：

\begin{itemize}
  \item \texttt{craftWoodenPickaxe}
  \item \texttt{craftStonePickaxe}
  \item \texttt{craftFurnace}
  \item \texttt{craftIronPickaxe}
  \item \texttt{mineThreeCobblestone}
  \item \texttt{mineCoalOre}
  \item \texttt{mineCopperOre}
  \item \texttt{mineLapisOre}
  \item \texttt{smeltFourRawCopper}
  \item \texttt{smeltFourRawIron}
\end{itemize}

\subsection{最后状态快照}

从 \texttt{voyager\_live.log} 尾部可见的最新状态如下：

\begin{itemize}
  \item Biome: \texttt{taiga}
  \item Position: \texttt{x=672.5, y=69.0, z=-399.5}
  \item 已装备物品：\texttt{iron\_boots}
  \item 关键库存：\texttt{iron\_pickaxe: 1}、\texttt{stone\_sword: 1}、\texttt{lapis\_lazuli: 6}、\texttt{raw\_iron: 5}、\texttt{copper\_ingot: 4}、\texttt{coal: 2}、\texttt{cobblestone: 66}
\end{itemize}

这说明本次运行已经稳定推进到铁工具与基础铁护甲阶段。

\section{终止原因分析}

根据现有日志，能够明确确认的最后一个内部失败点发生在任务 \texttt{Kill 1 skeleton} 期间。日志中出现了如下链条：

\begin{itemize}
  \item Python 侧通过 \texttt{voyager/env/bridge.py} 调用本地 mineflayer HTTP 服务 \texttt{/step}。
  \item 该请求在本地回环地址 \texttt{127.0.0.1:3000} 上触发了 \texttt{ReadTimeout}。
  \item 当时 \texttt{env\_request\_timeout} 已提升到 300 秒，说明这是一次长时间未返回的环境执行步骤，而不是短超时误判。
  \item 随后 Python 抛出 \texttt{RuntimeError: Failed to step Minecraft server}。
  \item Voyager 按设计将 \texttt{Kill 1 skeleton} 标记为失败，并继续进入下一轮 curriculum 问答流程。
\end{itemize}

因此，可以把 ``最后一个明确记录到的内部故障'' 定义为：

\begin{quote}
在 \texttt{Kill 1 skeleton} 任务中，本地 mineflayer \texttt{/step} 调用在 300 秒后超时，导致该任务失败。
\end{quote}

但是，日志尾部并没有留下另一个明确的 Python fatal traceback、\texttt{Interrupted by user}、或标准的 ``正常退出'' 标记。因此，``为什么整个进程最终停了'' 这一点无法从现有日志 100\% 精确复原。保守结论是：

\begin{itemize}
  \item 已知内部失败是 skeleton 任务的 \texttt{/step} 超时。
  \item 在这之后，程序仍继续输出了一段 curriculum QA 与状态快照。
  \item 整体停止更像是外部中断、会话结束、或日志未覆盖到的进程终止，而不是一个已经被完整记录的新代码异常。
\end{itemize}

\section{本次修复的效果}

从实际结果看，这些代码改动显著改善了运行稳定性：

\begin{itemize}
  \item 解决了早期 crafting/mineflayer 序列化问题，使木板、工作台、木镐等关键前置任务可以完成。
  \item 解决了自定义模型接口上 ``请求挂起但不报错'' 的问题，使 Action/Curriculum/Critic/Skill 四类 agent 都能继续推进。
  \item 解决了 mineflayer 生命周期中的空指针问题，降低 bot 重启过程中的二次崩溃概率。
  \item 通过提升环境超时与 \texttt{/start} 重试，减少了长时间地下/水下任务导致整轮学习直接退出的概率。
\end{itemize}

如果没有这些修复，本次运行基本不可能推进到铁镐与铁靴阶段。

\section{剩余问题与建议}

当前仍可见的问题主要有三类：

\begin{itemize}
  \item \textbf{战斗类任务稳定性不足}：\texttt{Kill 1 skeleton} 依然会在复杂环境下触发长时间 \texttt{/step} 卡住。
  \item \textbf{Critic fenced JSON}：critic 仍然经常先返回带代码围栏的 JSON，导致额外一次解析失败与重试。
  \item \textbf{mineflayer 运行时噪声}：日志中仍存在大量实体弃用 warning 和偶发连接断开提示，虽然不再总是致命，但会增加排障成本。
\end{itemize}

下一步建议：

\begin{enumerate}
  \item 优先增强战斗任务的动作超时与回退策略，特别是 skeleton 近战接近逻辑。
  \item 在 JSON 解析器中直接处理 fenced JSON，减少 critic 的无效重试。
  \item 对 \texttt{mineBlock} 和战斗任务分别增加更明确的成功/失败判据，避免 ``动作已结束但结果未达成'' 的模糊状态。
\end{enumerate}

\section{结论}

本次分支已经从 ``不稳定的论文代码运行环境'' 演化为 ``可在真实 Minecraft LAN 世界中持续推进的工程化版本''。虽然最终运行没有完整达到所有课程目标，但它已经真实完成 27 个任务、学习 28 个技能，并稳定推进到铁工具和基础铁护甲阶段。\texttt{Kill 1 skeleton} 仍是当前最明显的稳定性短板，但从现有代码与结果看，系统已经具备继续迭代和继续跑更长任务链的基础。

\end{document}
